\begin{initiativkarte}{3DSpace}

  % Bild (optional)
  \IfFileExists{bilder/3dspace.jpg}{%
    \begin{center}
      \includegraphics[width=\linewidth,height=5cm,keepaspectratio]{bilder/3dspace.jpg}
    \end{center}
    \vspace{0.4cm}
  }{}

  % Kurzbeschreibung
  \textit{\textcolor{hawgray}{Der 3DSpace ist ein offenes Labor der HAW Hamburg, in dem Studierende Ideen in reale Prototypen umsetzen und additive Fertigungsverfahren (3D-Druck) kennenlernen.}}

  \vspace{0.5cm}
  \hawrule

  % Beschreibung
  \textbf{\textcolor{hawblue}{Was machen wir?}}\\[0.3cm]

  Der 3DSpace der HAW Hamburg ist eine moderne Lern- und Entwicklungsumgebung für additive Fertigung und Prototyping.

  Studierende aller Fachrichtungen haben hier die Möglichkeit, eigene Ideen von der digitalen Konstruktion bis zum physischen Bauteil umzusetzen.

  Im Mittelpunkt steht die komplette digitale Prozesskette:
  von der CAD-Modellierung über die Fertigung bis hin zum fertigen Produkt.

  Zu den Aktivitäten gehören:
  \begin{itemize}
      \item Entwicklung und Umsetzung eigener Prototypen
      \item Nutzung von 3D-Druckern und additiven Fertigungsverfahren
      \item CAD-Modellierung und Design für Fertigung
      \item Arbeiten mit unterschiedlichen Materialien und Drucktechnologien
      \item Teilnahme an Workshops und betreuten Projekten
  \end{itemize}

  Der 3DSpace wird von Tutor:innen betreut und richtet sich sowohl an Einsteiger:innen als auch an fortgeschrittene Studierende, die praktische Erfahrung im Bereich moderner Fertigung sammeln möchten.

  \vspace{0.5cm}

  % Tags
  \tag{3D-Druck}
  \tag{Prototyping}
  \tag{CAD}
  \tag{Fertigung}
  \tag{Labor}

  \vspace{0.6cm}
  \hawrule

  % Infozeilen
  \infozeile{\faUsers}{Zielgruppe: Alle Studierenden der HAW Hamburg}
  \infozeile{\faGraduationCap}{Semester: Ab dem 1. Semester}
  \infozeile{\faMapMarker*}{Ort: HAW Hamburg, Berliner Tor 21 (Raum 035)}
  \infozeile{\faGlobe}{Website: \href{https://www.haw-hamburg.de/nachhaltige-ingenieurwissenschaften/laborzentren/3dspace/}{HAW 3DSpace}}

\end{initiativkarte}